% Retoca las líneas marcadas con TODO según las necesidades

\usepackage[spanish,es-tabla]{babel} % EMCR: la opcion es-table es para que en lugar de indice de cuadros aparezca indice de tablas
\usepackage[utf8]{inputenc}
\usepackage{geometry}
\usepackage{makeidx}
\usepackage{url}
\usepackage{graphicx}
\usepackage{caption}
\usepackage{acronym}
\usepackage{hyphenat}
\usepackage{a4wide}
\usepackage[normalsize]{subfigure}
\usepackage{float}
\usepackage{titlesec}
\usepackage[export]{adjustbox}
\usepackage[Lenny]{fncychap}
\usepackage{listings} % para poder hacer uso de "listings" propios (p.ej. códigos)
\usepackage{eurosym} % para poder usar el símbolo del euro con \euro {xx}
\usepackage{hyperref} % TODO: añade la opción hidelinks para imprimirlo (los enlaces no aparecerán resaltados)
\usepackage{graphicx}
\usepackage{float}
\usepackage{tikz}
\usepackage{array}
\usetikzlibrary{arrows.meta, positioning, shapes.geometric, fit, calc}
\usepackage{wrapfig}
\usepackage{tabularx}
\usepackage[table]{xcolor}
\usepackage{booktabs}
\usepackage{longtable}
\usepackage{pdflscape}

% Rojo  URJC
\definecolor{URJCRed}{HTML}{CB0017}
\definecolor{URJCRedDark}{HTML}{9E0012}
\definecolor{TableHeaderText}{HTML}{FFFFFF}
\definecolor{TableRowSoft}{HTML}{F7F7F7}
\definecolor{TableRule}{HTML}{B8B8B8}

% Colores para código estilo Arduino IDE en modo oscuro
\definecolor{ArduinoBg}{HTML}{1E1E1E}
\definecolor{ArduinoText}{HTML}{DCDCDC}
\definecolor{ArduinoComment}{HTML}{6A9955}
\definecolor{ArduinoKeyword}{HTML}{569CD6}
\definecolor{ArduinoFunction}{HTML}{DCDCAA}
\definecolor{ArduinoString}{HTML}{CE9178}
\definecolor{ArduinoNumber}{HTML}{B5CEA8}
\definecolor{ArduinoLine}{HTML}{858585}
\definecolor{ArduinoFrame}{HTML}{3C3C3C}

% Colores para código estilo bash linux en modo oscuro
\definecolor{TerminalBg}{HTML}{0C0C0C}
\definecolor{TerminalText}{HTML}{DCDCDC}
\definecolor{TerminalGreen}{HTML}{00C853}
\definecolor{TerminalBlue}{HTML}{569CD6}
\definecolor{TerminalCyan}{HTML}{4EC9B0}
\definecolor{TerminalYellow}{HTML}{DCDCAA}
\definecolor{TerminalOrange}{HTML}{CE9178}
\definecolor{TerminalGray}{HTML}{858585}
\definecolor{TerminalFrame}{HTML}{2D2D2D}

% Cabeceras de tabla: fila roja con texto blanco y negrita
\newcommand{\TVTableHeader}{%
    \rowcolor{URJCRed}%
    \color{TableHeaderText}\bfseries%
}

% Celda individual de cabecera, útil para asegurar color blanco en todas las columnas
\newcommand{\TVHead}[1]{\textcolor{TableHeaderText}{\textbf{#1}}}

% Estilo de captions de tablas
\captionsetup[table]{
    labelfont={bf,color=URJCRed},
    textfont=it,
    justification=centering,
    singlelinecheck=false,
    skip=6pt
}

% Estilo de captions de códigos/listados
\captionsetup[lstlisting]{
    labelfont={bf,color=URJCRed},
    textfont=it,
    justification=centering,
    singlelinecheck=false,
    skip=6pt
}

% Lenguaje Arduino para listings

\lstdefinelanguage{Arduino}{
    language=C++,
    morekeywords={
        setup,loop,pinMode,digitalWrite,digitalRead,analogRead,analogWrite,
        delay,delayMicroseconds,millis,micros,Serial,SoftwareSerial,Wire,
        HIGH,LOW,INPUT,OUTPUT,INPUT_PULLUP,true,false,
        String,byte,boolean,unsigned,long,uint8_t,uint16_t,uint32_t
    },
    sensitive=true,
    morecomment=[l]{//},
    morecomment=[s]{/*}{*/},
    morestring=[b]"
}

% Estilo de código oscuro tipo Arduino IDE

\lstdefinestyle{arduinoDark}{
    backgroundcolor=\color{ArduinoBg},
    basicstyle=\small\ttfamily\color{ArduinoText},
    keywordstyle=\color{ArduinoKeyword}\bfseries,
    commentstyle=\color{ArduinoComment}\itshape,
    stringstyle=\color{ArduinoString},
    numberstyle=\tiny\color{ArduinoLine},
    identifierstyle=\color{ArduinoText},
    frame=single,
    rulecolor=\color{ArduinoFrame},
    framerule=0.4pt,
    numbers=left,
    numbersep=8pt,
    showstringspaces=false,
    columns=fullflexible,
    keepspaces=true,
    breaklines=true,
    breakatwhitespace=true,
    tabsize=2,
    captionpos=b,
    xleftmargin=1.2em,
    xrightmargin=0.5em,
    framexleftmargin=1em,
    aboveskip=3mm,
    belowskip=3mm
}

% Lenguajes y estilos para terminal Raspberry Pi / Linux

\lstdefinelanguage{RaspberryBash}{
    morekeywords={
        sudo,apt,update,upgrade,install,git,clone,cd,nano,chmod,curl,
        mkdir,rsync,python3,pip,source,uvicorn,systemctl,tail,cat,grep,
        echo,export,pull,start,stop,restart,status,logs
    },
    sensitive=true,
    morecomment=[l]{\#},
    morestring=[b]",
    morestring=[b]'
}

\lstdefinelanguage{TerminalPlain}{}

\lstdefinelanguage{JsonTerminal}{
    morestring=[b]",
    morekeywords={true,false,null},
    sensitive=true
}

\lstdefinestyle{raspberryTerminal}{
    language=RaspberryBash,
    backgroundcolor=\color{TerminalBg},
    basicstyle=\footnotesize\ttfamily\color{TerminalText},
    keywordstyle=\color{TerminalGreen}\bfseries,
    commentstyle=\color{TerminalGray}\itshape,
    stringstyle=\color{TerminalOrange},
    emph={
        tierraviva.sh,main.py,api.py,database.py,index.html,
        install,start,stop,restart,status,logs,
        dashboard,frontend,images
    },
    emphstyle=\color{TerminalCyan}\bfseries,
    frame=single,
    rulecolor=\color{TerminalFrame},
    framerule=0.5pt,
    numbers=none,
    showstringspaces=false,
    columns=fullflexible,
    keepspaces=true,
    breaklines=true,
    breakatwhitespace=true,
    tabsize=2,
    captionpos=b,
    xleftmargin=0.6em,
    xrightmargin=0.4em,
    framexleftmargin=0.6em,
    aboveskip=2.5mm,
    belowskip=2.5mm
}

\lstdefinestyle{raspberryOutput}{
    language=TerminalPlain,
    backgroundcolor=\color{TerminalBg},
    basicstyle=\footnotesize\ttfamily\color{TerminalText},
    frame=single,
    rulecolor=\color{TerminalFrame},
    framerule=0.5pt,
    numbers=none,
    showstringspaces=false,
    columns=fullflexible,
    keepspaces=true,
    breaklines=true,
    breakatwhitespace=true,
    tabsize=2,
    xleftmargin=0.6em,
    xrightmargin=0.4em,
    framexleftmargin=0.6em,
    aboveskip=2.5mm,
    belowskip=2.5mm
}

\lstdefinestyle{jsonTerminal}{
    language=JsonTerminal,
    backgroundcolor=\color{TerminalBg},
    basicstyle=\footnotesize\ttfamily\color{TerminalText},
    keywordstyle=\color{TerminalBlue}\bfseries,
    stringstyle=\color{TerminalOrange},
    frame=single,
    rulecolor=\color{TerminalFrame},
    framerule=0.5pt,
    numbers=none,
    showstringspaces=false,
    columns=fullflexible,
    keepspaces=true,
    breaklines=true,
    breakatwhitespace=true,
    tabsize=2,
    xleftmargin=0.6em,
    xrightmargin=0.4em,
    framexleftmargin=0.6em,
    aboveskip=2.5mm,
    belowskip=2.5mm
}

\lstnewenvironment{bashcode}[1][]
{\lstset{style=raspberryTerminal,#1}}
{}

\lstnewenvironment{terminaloutput}[1][]
{\lstset{style=raspberryOutput,#1}}
{}

\lstnewenvironment{jsoncode}[1][]
{\lstset{style=jsonTerminal,#1}}
{}

% Estilo por defecto para todos los lstlisting
\lstset{
    style=arduinoDark,
    language=Arduino
}

\newcommand{\imgmaterial}[2]{%
    \IfFileExists{#1}{%
        \includegraphics[width=#2,keepaspectratio]{#1}%
    }{%
        \fbox{\parbox[c][1.4cm][c]{#2}{\centering\scriptsize Imagen\\no encontrada}}%
    }%
}
% Para que no parta las palabras
\pretolerance=10000

\newcommand{\bigrule}{\titlerule[0.5mm]} \titleformat{\chapter}[display] % cambiamos el formato de los capítulos
{\bfseries\Huge} % por defecto se usaron caracteres de tamaño huge en negrita
{% contenido de la etiqueta 
\titlerule % línea horizontal 
\filright % texto alineado a la derecha 
\Large\chaptertitlename\ % capítulo e índice en tamaño large
\Large % en lugar de 
\Huge \Large\thechapter} 
{0mm} % espacio mínimo entre etiqueta y cuerpo
{\filright} % texto del cuerpo alineado a la derecha
[\vspace{0.5mm} \bigrule] % después del cuerpo, dejar espacio vertical y trazar línea horizontal gruesa
\geometry{a4paper, left=3.5cm, right=2cm, top=3cm, bottom=2cm, headsep=1.5cm}

% Estilos para ilustrar códigos:
\definecolor{code_green}{rgb}{0,0.6,0}
\definecolor{code_gray}{rgb}{0.5,0.5,0.5}
\definecolor{code_mauve}{rgb}{0.58,0,0.82}


% Definición de mis propios tipos: Códigos, Ecuaciones y Tablas
\DeclareCaptionType{code}[Código][Listado de códigos]
\DeclareCaptionType{myequation}[Ecuación][Listado de ecuaciones]

% TODO: especifica las reglas de separación que consideres. Algunos ejemplos:
\hyphenation{fuer-tes}
\hyphenation{mul-ti-ca-pa}
\hyphenation{res-pues-ta}
\hyphenation{di-fe-ren-tes}
\hyphenation{de-sa-rro-lla-dos}
\hyphenation{re-pre-sen-tan-do}

